\documentclass{article}
\usepackage{etoolbox}
\usepackage{graphicx}

\ifdef{\myoutput}{}{\def\myoutput{concept}}
\ifdef{\myversion}{}{\def\myversion{1}}

\usepackage[output=\myoutput
,numberofversions=4
,version=\myversion
,seed=4
,randomizequestions=true
,randomizeanswers=true
,writeRfile=true
]{mcexam}

\usepackage[margin=0.95in]{geometry}
\usepackage{multicol}

\usepackage{fancyhdr,lastpage}
\pagestyle{fancy}
\fancyhf{}
\renewcommand{\headrulewidth}{1pt} 
\renewcommand{\footrulewidth}{1pt}
\lfoot{\mctheversion}
\rfoot{Page \thepage\ of \pageref{LastPage}}
\lhead{Quiz 7.3 and 8.1 \mctheversion \\ MAT 2000: Introduction to Statistics}
\rhead{Name Print): \hspace{2.5 in}}

\usepackage{calc}
\renewenvironment{setmcquestion}
{\begin{minipage}[t]{\linewidth-\labelwidth}}
	{\end{minipage}\par}
	
\renewenvironment{setmcquestioninstruction}
{\begin{minipage}{\textwidth}}
	{\end{minipage}}
	
\renewenvironment{setmcanswers}{}{}
\setlist[setmcquestions]{label=\mcquestionlabelfmt{*}.
	,ref=\mcquestionlabelfmt{*}
	,itemsep=0.5\baselineskip
	,topsep=1\baselineskip
}

\renewcommand\mcversionlabelfmt[1]{\arabic{#1}}

%copy-pasting from MS Word, use regex to find [\*]*[a-d]\. and replace with \answer 

\raggedbottom

\begin{document}
	Circle the correct answer for each question.  If you're unsure, you can write numbers 1, 2, 3, etc. next to the choices to get partial credit.  Questions refer to the nonindented paragraph above unless otherwise stated.\\[.5cm]
	
	\begin{mcquestions}
			
		\begin{mcquestioninstruction}
			A random sample of Hope College students was taken and one of the questions asked was how many hours per week they study.  The goal is to see whether a difference exists between males and females in terms of average study time.  The sample results are given in the following table.
			\begin{center}
				\begin{tabular}{| l | l | l |}
					\hline
					& \textbf{Female} & \textbf{Male} \\
					\hline
					\textbf{Sample Size} & 125 & 75 \\
					\hline
					\textbf{Sample Average} & 18.52 & 15.17 \\
					\hline
					\textbf{Sample SD} & 10.85 & 11.32 \\
					\hline
				\end{tabular}
			\end{center}
		\end{mcquestioninstruction}
		
		\question Assuming the distribution of study times is not strongly skewed for either sample, which approach would be more appropriate for these data: simulation-based or theory-based?
		\begin{mcanswerslist}
			\answer Simulation-based, since this approach works for any sample size.
			\answer Theory-based, since both 125 and 75 are greater than 20.
			\answer[correct] Either approach would be appropriate.
			\answer Neither approach would be appropriate.
		\end{mcanswerslist}
		
		\begin{mcquestioninstruction}
			An article that appeared in the British Medical Journal (2010) presented the results of a randomized experiment conducted by researcher Jeremy Groves, whose objective was to determine whether the weight of his bicycle could affect his travel time to work. On each of 56 days (from mid-January to mid-July 2010), Groves tossed a £1 coin to decide whether to bike to work on a carbon frame (lighter) bicycle that weighed 20.9 lbs or a steel frame (heavier) bicycle that weighed 29.75 lbs, and recorded the commute time (in minutes) for each trip.
			
			Here are the summary statistics for his data:
			
			\begin{center}
				\begin{tabular}{| l | l | l | l |}
					\hline
					\textbf{Bike Type} & \textbf{Sample size} & \textbf{Sample average} & \textbf{Sample SD} \\
					\hline
					\textbf{Carbon frame (lighter)} & 26 & 108.34 min & 6.25 min \\
					\hline
					\textbf{Steel frame (heavier)} & 30 & 107.81 min & 3.89 min \\
					\hline
				\end{tabular}
			\end{center}
		\end{mcquestioninstruction}
		
		\question Calculate the standardized statistic for the appropriate test of two means (carbon – steel).
		\begin{mcanswerslist}
			\answer 0.71
			\answer 0.53
			\answer[correct] 0.37
			\answer 1.42
		\end{mcanswerslist}
		
		\question[follow] Under the null hypothesis, what distribution does the test statistic follow?
		\begin{mcanswerslist}
			\answer Standard normal distribution
			\answer[correct] t-distribution
			\answer Skewed distribution
			\answer Chi-squared distribution
		\end{mcanswerslist}
		
		\begin{mcquestioninstruction}
			In order to investigate whether talking on cell phones is more distracting than listening to car radios while driving, sixty-four student volunteers (from a single college class) were randomly assigned to a cell phone group or a radio group (32 students were assigned to each group). Each student ``drove'' a machine that simulated driving situations. While ``driving'' the simulator, a target would flash red at irregular intervals. Participants were instructed to press the ``brake'' button as soon as possible when they detected a red light. Participant response times were measured as the time between the red light appearing and pushing the brake button. While driving, the radio group listened to a radio broadcast and the cell phone group carried on a conversation on the cell phone with someone in the next room.
		\end{mcquestioninstruction}
		
		\question The theory based 99\% confidence interval for $\mu_{{cell}}-\mu_{{radio}}$ is given by $(-0.7401,103.7401)$. Does the interval provide significant statistical evidence that the long-run mean response time differs between the cell phone and radio treatments?
		\begin{mcanswerslist}
			\answer Yes, since most of the interval is positive.
			\answer Yes, since zero is not contained in the interval.
			\answer[correct] No, since zero is contained in the interval.
			\answer The interval does not provide enough information to answer this question.
		\end{mcanswerslist}
		
		\question A 1997 study described in the European Journal of Clinical Nutrition compares the growth of a random sample of vegetarian and omnivorous children, ages 7 - 11, in Northwest England. One of the aspects on which the children were compared was their body mass index (BMI). Of the 50 vegetarian children in the study, the average BMI was 16.76 and the standard deviation was 1.91. Among the 30 omnivorous children in the study, the average BMI was 17.12 and the standard deviation was 2.23.
		
		Which best describes this study design?
		\begin{mcanswerslist}
			\answer Paired using repeated measures
			\answer Paired using matching
			\answer[correct] Independent groups
			\answer None of the above
		\end{mcanswerslist}
		
		\question A random sample of medical records from 28 heart attack patients at a Boston hospital was used to better understand the relationship between cholesterol and heart failure.  The cholesterol levels of heart attack patients were recorded at 2 days and at 14 days after the patient had a heart attack.
		
		Which best describes this study design?
		\begin{mcanswerslist}
			\answer[correct] Paired using repeated measures
			\answer Paired using matching
			\answer Independent groups
			\answer None of the above
		\end{mcanswerslist}
		
		\question Researchers randomly selected 36 children who had been identified as gifted from a large city. The IQ of each child and the IQ of the child's mother were measured. A statistical analysis will be performed to determine if IQs of gifted children tend to be higher than their mothers' IQ on average.
		
		Which best describes this study design?
		
		\begin{mcanswerslist}
			\answer Paired using repeated measures
			\answer[correct] Paired using matching
			\answer Independent groups
			\answer None of the above
		\end{mcanswerslist}
	\end{mcquestions}
\end{document}